\begin{ouverture}[grandeurs]{11}{Longueurs et périmètres}{Partie 3 · Grandeurs et mesures}
{Unités · Conversions raisonnées · Tour d’une figure}
\objectif{Choisir une unité de longueur et convertir en raisonnant.}
\objectif{Estimer une longueur avec des références connues.}
\objectif{Calculer le périmètre d’un polygone.}
\prerequis{Nombres décimaux et multiplication par 10 ou 100.}
\end{ouverture}
\pagestyle{otMain}\essentiel\fontsize{11}{15}\selectfont
\phantomsection\label{cours-11-1}
\essSection{1}{Mesurer des longueurs}
\begin{definition}Les unités de longueur vont du millimètre au kilomètre : mm, cm, dm, m, km. Relations à connaître : $1\ \mathrm{cm}=10\ \mathrm{mm}$, $1\ \mathrm m=100\ \mathrm{cm}$, $1\ \mathrm m=1\,000\ \mathrm{mm}$ et $1\ \mathrm{km}=1\,000\ \mathrm m$.\end{definition}
\begin{savoirfaire}[Convertir en raisonnant]Convertir en s’appuyant sur une relation : $3{,}5$ m, c’est $350$ cm, car 1 m vaut 100 cm et $3{,}5\times100=350$. De même, $48$ mm $=4{,}8$ cm, car 10 mm font 1 cm.\end{savoirfaire}
\begin{exemple}Des références aident à estimer : une porte mesure environ 2 m, la largeur d’un doigt environ 1 cm, l’épaisseur d’une pièce de monnaie environ 2 mm, un tour de stade 400 m.\end{exemple}
\begin{attention}Avant de comparer ou d’additionner des longueurs, les écrire dans la même unité : $0{,}5$ m et 48 cm se comparent en centimètres, $50>48$.\end{attention}
\clearpage
\phantomsection\label{cours-11-2}
\essSection{2}{Le périmètre d’une figure}
\begin{definition}Le périmètre d’une figure est la longueur de son contour. Pour un polygone, on additionne les longueurs de tous les côtés, données ou mesurées à la règle graduée.\end{definition}
\label{fig-cours-11-2-cours-perimetre}\figureGeo{cours-perimetre}

\begin{savoirfaire}[Des règles que l’on retrouve]Le carré a quatre côtés de même longueur : son périmètre est le quadruple du côté. Le rectangle a deux longueurs et deux largeurs : on peut calculer la somme longueur plus largeur, puis la doubler.\end{savoirfaire}
\begin{exemple}Un rectangle de 8 cm sur $3{,}5$ cm a pour périmètre $8+3{,}5+8+3{,}5=23$ cm. Pour une figure en L, repérer tous les côtés : certains se déduisent des autres avant d’additionner.\end{exemple}

\clearpage\rubrique{Découverte}
\begin{blocLivre}{ex-11-1}\begin{exercice}[1]{La bonne unité}
Choisir l’unité adaptée (mm, cm, m ou km) : épaisseur d’une pièce de monnaie ; longueur d’un cahier ; hauteur d’un immeuble ; distance entre deux villes.
\end{exercice}\end{blocLivre}
\begin{blocLivre}{ex-11-2}\begin{exercice}[1]{Convertir en raisonnant}
Convertir : $3{,}5$ m en cm ; 2 km en m ; 48 mm en cm.
\end{exercice}\end{blocLivre}
\begin{blocLivre}{ex-11-3}\begin{exercice}[1]{Comparer des longueurs}
Comparer $0{,}5$ m et 48 cm, puis $1\,200$ m et $1{,}5$ km.
\end{exercice}\end{blocLivre}
\begin{blocLivre}{ex-11-4}\begin{exercice}[1]{Estimer sans mesurer}
Associer chaque objet à une longueur : une gomme ; un terrain de football ; la hauteur d’une chaise. Longueurs proposées : 100 m ; 4 cm ; 90 cm.
\end{exercice}\end{blocLivre}
\begin{blocLivre}{ex-11-5}\begin{exercice}[1]{Le tour du quadrilatère}
Calculer le périmètre du quadrilatère dont les côtés mesurent 3 cm, $4{,}5$ cm, $2{,}5$ cm et 4 cm.
\label{fig-ex-11-5}\figureGeo{per-5}
\end{exercice}\end{blocLivre}

\rubrique{Entraînement}
\begin{blocLivre}{ex-11-6}\begin{exercice}[2]{Le carré}
Un carré a des côtés de 7 cm. Calculer son périmètre.
\label{fig-ex-11-6}\figureGeo{per-6}
\end{exercice}\end{blocLivre}
\begin{blocLivre}{ex-11-7}\begin{exercice}[2]{Le rectangle}
Calculer le périmètre d’un rectangle de longueur 8 cm et de largeur $3{,}5$ cm.
\label{fig-ex-11-7}\figureGeo{per-7}
\end{exercice}\end{blocLivre}
\begin{blocLivre}{ex-11-8}\begin{exercice}[2]{Retrouver le côté}
Le périmètre d’un carré vaut 36 cm. Quelle est la longueur d’un côté ?
\end{exercice}\end{blocLivre}
\begin{blocLivre}{ex-11-9}\begin{exercice}[2]{Des kilomètres aux mètres}
Convertir $1{,}25$ km en mètres, puis 230 cm en mètres.
\end{exercice}\end{blocLivre}
\begin{blocLivre}{ex-11-10}\begin{exercice}[2]{La figure en L}
Calculer le périmètre de la figure en L, dont les mesures sont indiquées sur le dessin.
\label{fig-ex-11-10}\figureGeo{per-10}
\end{exercice}\end{blocLivre}

\rubrique{Problèmes}
\begin{blocLivre}{ex-11-11}\begin{exercice}[3]{La clôture du jardin}
Un jardin rectangulaire mesure $12{,}5$ m sur 8 m. On veut le clôturer en laissant un portail de $1{,}5$ m.\begin{questions}
\item Calculer le périmètre du jardin.
\item Quelle longueur de grillage faut-il acheter ?
\end{questions}
\end{exercice}\end{blocLivre}
\begin{blocLivre}{ex-11-12}\begin{exercice}[3]{Les tours de piste}
Une piste d’athlétisme mesure 400 m par tour. Un coureur effectue 5 tours. Quelle distance parcourt-il, en mètres puis en kilomètres ?
\end{exercice}\end{blocLivre}
\begin{blocLivre}{ex-11-13}\begin{exercice}[3]{Deux chemins sur quadrillage}
Sur le quadrillage, chaque carreau mesure 1 cm de côté. Comparer les périmètres du rectangle et de la figure en escalier.
\label{fig-ex-11-13}\figureGeo{per-13}
\end{exercice}\end{blocLivre}
\begin{blocLivre}{ex-11-14}\begin{exercice}[3]{Le cadre de la photo}
Une photo rectangulaire mesure 15 cm sur 10 cm. On veut coller un ruban sur tout son contour. Le ruban disponible mesure $0{,}45$ m. Suffira-t-il ?
\end{exercice}\end{blocLivre}
\begin{blocLivre}{ex-11-15}\begin{exercice}[3]{Même périmètre, même figure ?}
Lou affirme : « Deux figures qui ont le même périmètre sont forcément identiques. » Trouver un contre-exemple avec un carré et un rectangle.
\end{exercice}\end{blocLivre}

\clearpage
